\begin{abstract}

Cloud-native systems increasingly rely on infrastructure services (e.g., service meshes, monitoring agents), which compete for resources with user applications, degrading performance and scalability.
We propose \xpod, a new abstraction that offloads these services to Data Processing Units (DPUs) to enforce strict isolation while reducing host resource contention and operational costs.
To realize \xpod, we introduce \os, a cross-PU (XPU) network system featuring: (1) \emph{split network namespace}, a unified network abstraction for processes spanning CPU and DPU,
and (2) \emph{elastic and efficient XPU networking}, a communication mechanism achieving shared-memory performance without pinned resource overhead and polling costs.
%These innovations enable \emph{dynamic split} --- runtime decoupling of multi-container Pods or multi-process containers across heterogeneous PUs, requiring no application modifications.
By leveraging \os and the compositional nature of cloud-native workloads, \xpod can optimally offload infrastructure containers to DPUs.
We implement \os based on Linux, and implement a cloud-native system called \sys based on Kubernetes.
We evaluate the systems using NVIDIA Bluefield-2 DPUs and CXL-based DPUs (simulated with real CXL memory devices).
The results show that \sys effectively supports complex (unmodified) commodity cloud-native applications (up to 1 million LoC) and provides up to 31.9x better latency and 64x less resource consumption (compared with kernel-bypass design), 60\% better end-to-end latency, and 55\% higher scalability compared with SOTA systems.

\end{abstract}

